% !TeX program = pdflatex
\documentclass[11pt,a4paper]{article}

% --- Język i czcionki ---
\usepackage[T1]{fontenc}
\usepackage[utf8]{inputenc}
\usepackage[polish]{babel}
\usepackage{lmodern}
\usepackage[final]{microtype}

% --- Układ strony i nagłówki/stopki ---
\usepackage[a4paper,margin=2.5cm]{geometry}
\usepackage{fancyhdr}
\pagestyle{fancy}
\fancyhf{} % czyść wszystko
\lhead{Etap 2.}
\rhead{\leftmark}
\cfoot{\thepage}

% --- Matematyka ---
\usepackage{amsmath,amssymb,amsthm,mathtools}
\usepackage{siunitx} % jeżeli potrzebne jednostki
\sisetup{locale = PL} % separator 1,23 (polski), można zmienić na locale=PL gdy dostępne
\usepackage{bm}       % pogrubione symbole

% --- Odsyłacze i tytuły ---
\usepackage[hidelinks]{hyperref}
\usepackage[nameinlink,capitalise,noabbrev]{cleveref}
\usepackage{csquotes}

% --- Listy i drobiazgi ---
\usepackage{enumitem}
\setlist{noitemsep,topsep=3pt}
\usepackage{xcolor}

% --- Numeracja i wygląd sekcji ---
\numberwithin{equation}{section}
\setcounter{secnumdepth}{3}
\setcounter{tocdepth}{2}

% --- Meta ---
\title{Sprawozdanie z etapu 3.}
\author{Jakub Kogut}
\date{\today}

\begin{document}
\maketitle
W tym etapie projektu przedstawiamy nasz plan na jego implementację. Składa się ona z odpowiednio z projektu backendu oraz frontendu.
\section{Backend}
Nasz backend będzie odpowiedziany za przechowywanie oraz przetwarzanie wszelakich danych jak i musi posiadać \textit{user friendly} metodę na ich wpowadzenie. W tym celu zdecydowaliśmy się na użycie frameworku \textbf{Django} wraz z \textbf{Django REST Framework} do stworzenia API. Do przechowywania danych użyjemy bazy danych \textbf{PostgreSQL}, która jest utrzymywana przez \textbf{Django}.
\subsection{Baza danych}
W bazie danych przechowywane będą wszelkie dane dotyczące użytkowników jak i zdobytych przez nich nagród, czy dane do kursów oraz quizów, a także forum.
\subsubsection{Tabele}
W bazie danych znajdować się będą następujące tabele:
\begin{itemize}
    \item \textbf{Users} - tabela przechowująca dane użytkowników
    \item \textbf{Courses} - tabela przechowująca dane kursów
\end{itemize}


\subsubsection{Funkcje, Triggery}

\subsection{Permisje}
Aby zabezpieczyć aplikację przed nieautoryzowanym dostępem, wprowadzimy system permisji. Użytkownicy będą podzieleni na dwie grupy:
\begin{itemize}
    \item \textbf{Użytkownicy zarejestrowani}\footnote{restracja poprzez nickname i hasło} - będą mieli dostęp do kursów, quizów oraz forum, nagród innych użytkowników
    \item \textbf{Administratorzy} - będą mieli dostęp do interfejsu administratora, gdzie będą mogli zarządzać kursami, quizami oraz użytkownikami
\end{itemize}
Dodatkowo, jeżeli użytkownik chciałby edytować/usunąć swój post/komentarz na forum, będzie wymagane sprawdzenie czy jest on autorem danego postu/komentarza.



\section{Interfejs administratora}
Django posiada wbudowany interfejs administratora, który pozwala na zarządzanie bazą danych bez konieczności tworzenia własnego panelu. W naszym projekcie wykorzystamy ten interfejs do zarządzania w szególności kursami oraz quizami.\\
Administrator będzie miał możliwość tworzenia oraz edytowania poprzez interfejs administratora:
\begin{itemize}
    \item tworzenie lekcji na podstawie dodawania slaidów z tekstem, ewentualnie obrazkami
    \item tworzenie quizów na podstawie pytań otwartych (warości numeryczne) oraz zamkniętych.
    \item dodawanie avatarów użytkowników (userzy mogą wybierać avatary z puli dostępnych w serwisie)
\end{itemize}
\end{document}

